%% GENERATED by python/paper/build.py from results/ -- do not edit by hand.
%% Rebuild:  .venv\Scripts\python.exe python\paper\build.py --only OECD26_Layer1_omegaUS
%% Source:   results/oecd26/layer1.csv, layer2.csv, anchor.csv; data/oecdFigure1.csv, oecd26_theta.csv
\textit{Note:} 21 countries in the data and in the model, with the least-squares lines whose slopes table~\ref{table:OECD26:slopes_omegaUS} reports. The model\textquotesingle s index is $1/(2-\theta)$ of the design in force in 2020, chosen in 1990, and its spending $\tau(1-\alpha)$ in 2020 on the path of the design it chooses. The model: every country at the U.S.\ $\beta = 0.65$, $\omega = 1.26$ and $\kappa = 97.29$, each at its cost $\lambda_{n,t} = 2\kappa\,\pi_{n,t}(1-\pi_{n,t})$, $\pi_{n,t} = \Pi_{n,t}/R_{n,t}$. Turnout: each country\textquotesingle s own propensities to vote by income group, from the CSES; Luxembourg, which has no CSES survey and enforces compulsory voting, at a flat profile.
